% ================================================================
%  chapters/parts/ch12-a.tex
%  第 12 章  扫描探针显微镜与纳米光谱技术（分片 a）
%  源：S2（下册）第 12 章，书内 p.885—912 / PDF p.372—398
%  范围：12.1（STM 基本原理：结构、隧穿电流、STS）、12.2（STM 在催化研究中的应用）
%  不含：12.3 及以后（AFM、TERS、Nano-FTIR，见分片 b/c）
% ================================================================

\chapterstart{扫描探针显微镜与纳米光谱技术}%
{本章回答：原子尺度的针尖凭什么能把表面看清楚，看到的究竟是几何高度还是电子结构，以及在
什么条件下这些结果才能外推到真实催化体系。全章覆盖扫描隧道显微镜（STM）、原子力显微镜
（AFM）与两类探针增强纳米光谱；分片 a 覆盖 12.1 成像原理与 12.2 催化应用。}
\label{ch:12}

\subsection{扫描隧道显微镜（STM）基本原理}
\label{sec:stm-principle}

\subsubsection{STM 的基本工作原理与结构}

\begin{probout}
扫描隧道显微镜（scanning tunneling microscopy，STM）把一个\textbf{量子隧穿结}当作位移
传感器：金属针尖与导电样品在真空中相距约 \qty{1}{nm} 或更近，加上偏压后产生隧穿电流；
针尖在 $xy$ 平面内逐点扫描，同时记录电流（或为维持电流所调整的 $z$ 位移），即得表面
图像。原始输出只有一组 $I(x,y)$ 或 $z(x,y)$ 数据，它同时编码了\textbf{几何起伏}与
\textbf{电子态密度}两种信息，二者在原理上不可分离。
\end{probout}

STM 的仪器构造分为机械设计与电子线路两部分：机械部分包括扫描管、粗调定位器、针尖、样品台
与振动隔离系统；电子部分包括电流放大器、反馈控制系统与扫描控制系统\src{S2 §12.1.1，
书内 p.888—889 / PDF p.374—375}。扫描管与粗调定位器都由压电陶瓷构成：管式扫描器的外管壁
镀四个电极分别对应 $\pm x$ 与 $\pm y$ 象限，在相对电极上加相反电压即可控制针尖在 $xy$
平面内的运动；管内壁接 $z$ 电极，控制垂直方向的伸缩。压电陶瓷的伸缩量与所加电压有精确的
线性关系（逆压电效应），压电系数在 \qty{}{\angstrom/V} 量级，因此针尖位置可精确控制到
\qty{0.1}{\angstrom}。

隧穿电流信号极小，典型值仅 \qtyrange{0.01}{10}{nA}\src{S2 §12.1.1，书内 p.889 /
PDF p.375}，须先由前级放大器放大并转换为电压。反馈系统计算实测电流与设定电流（参考值）的
误差并放大输出到 $z$ 压电元件：电流大于设定值则使针尖回撤远离样品，反之进给。这一负反馈
机制给出针尖-样品的平衡位置；针尖沿 $xy$ 扫描时，其 $z$ 向平衡位置描绘的就是等隧穿电流
轮廓，即恒流模式 STM 图像。振动隔离同样关键：除增加扫描单元刚性外，一般把扫描单元悬挂在
减震弹簧上并配涡流阻尼减震器，把环境振动影响减小至 \qty{1}{pm} 以下——整台仪器是一支
量程不到 \qty{1}{nm}、灵敏度却进入皮米量级的位移传感器，任何机械扰动都会直接进入信号。

\begin{figure}[htbp]
\centering
\begin{tikzpicture}[
  x=1cm, y=1cm,
  axisln/.style={thinline, black},
  guide/.style={thinline, black, dashed},
  ann/.style={font=\fontsize{7.2}{8.6}\selectfont, align=center},
  annl/.style={ann, align=left},
  fml/.style={font=\fontsize{8}{10}\selectfont}
]
% ---------- 样品 ----------
\fill[black!12] (6.0,0.05) rectangle (8.2,0.55);
\draw[mainline] (6.0,0.55) -- (8.2,0.55);
\foreach \xx in {6.25,6.75,7.25,7.75} {
  \fill[black!30] (\xx,0.55) circle (1.7pt);
  \draw[thinline] (\xx,0.55) circle (1.7pt);
}
\node[ann] at (7.1,-0.10) {样品};
% ---------- 吸附质 ----------
\fill[black!65] (6.42,0.98) circle (2.3pt);
\node[ann] at (5.55,1.30) {吸附质};
\draw[thinline] (5.85,1.24) -- (6.36,1.02);
\draw[guide] (6.42,0.60) -- (6.42,1.16);
% ---------- 针尖 ----------
\fill[black!22]
  (1.20,3.00) -- (1.40,1.55) .. controls (1.95,1.24) and (2.25,1.24) ..
  (2.80,1.55) -- (3.00,3.00) -- cycle;
\draw[mainline]
  (1.20,3.00) -- (1.40,1.55) .. controls (1.95,1.24) and (2.25,1.24) ..
  (2.80,1.55) -- (3.00,3.00);
\node[ann] at (2.10,2.62) {金属针尖};
% ---------- 势垒区 ----------
\fill[black!8] (3.00,0.55) rectangle (6.00,3.00);
% 【重叠修正，2026-09-30】`势垒区（真空）` 原用 anchor=north west 挂在 (3.08,2.12)，
% 两行基线在 1.91 / 1.60 cm，而 `功函数 $\Phi$` 挂在 (3.32,1.60)（anchor=west，
% 两行基线 1.69 / 1.38），实测前者第二行与后者第一行只差 0.09 cm——
% 7 对字形墨迹相交（最大 `功`×`垒` 0.543 pt²）。改挂势垒区右上角 anchor=north east。
\node[annl, anchor=north east] at (5.95,2.90) {势垒区\\（真空）};
\draw[axisln] (3.00,1.95) -- (6.00,1.95);
\node[ann, anchor=west] at (6.06,1.95) {真空能级};
\draw[axisln] (3.00,0.80) -- (6.00,0.80);
\node[ann, anchor=west] at (5.20,0.66) {$E_{\mathrm{F}}$};
\draw[guide] (3.10,1.95) -- (3.02,0.80);
\node[ann, align=left, anchor=west] at (3.32,1.60) {功函数\\ $\Phi$};
% ---------- 波函数在势垒中的衰减 ----------
\draw[mainline, dashed]
  plot[domain=3.02:5.98, samples=12, variable=\t]
  ({\t},{1.10+1.45*exp(-2.05*(5.98-\t))});
\node[ann, align=left, anchor=west] at (3.10,2.52)
  {电子波函数\\ 向真空衰减};
\node[fml] at (4.55,-0.12) {$I\propto V\mathrm{e}^{-2\kappa d}$};
% ---------- 间距标注 ----------
\draw[guide] (2.90,1.40) -- (6.10,1.40);
\draw[guide] (3.00,1.40) -- (3.00,1.02);
\draw[guide] (6.00,1.40) -- (6.00,1.02);
% 【重叠修正，2026-09-30】`$d$：针尖-样品间距` 原居中挂在 (4.50,1.02)，右端伸到
% x≈5.36 cm，压住 $E_{\mathrm{F}}$（挂在 5.20）：实测 `距`×`E` 墨迹相交 4.6×2.9 pt、
% 轮廓交叠 0.883 pt²。左移 0.25 cm，右端止于 5.11 cm，与 E 标签留 0.09 cm。
\node[ann, anchor=north] at (4.25,1.02) {$d$：针尖-样品间距};
% ---------- 样品偏压 ----------
\draw[axisln] (6.60,0.05) -- (6.60,-0.62);
\draw[axisln] (6.36,-0.62) -- (6.84,-0.62);
\node[ann, anchor=west] at (6.90,-0.62) {$V$};
\end{tikzpicture}
\caption{针尖-样品隧穿结与势垒的能级示意。读者应看两处：电子波函数在势垒区
中的指数衰减、功函数 $\Phi$ 决定的势垒高度（$E_{\mathrm{F}}=-\Phi$）。恒流反馈
回路（隧道电流放大器 $\to$ 反馈控制 $\to z$ 压电元件）与偏压下的能级对齐过程
见正文 §12.1.1--§12.1.2，本图只画隧穿结。图中横向尺寸与实际比例不符，为示意
画法。\schemamark\figsrc{据 S2（下册）图 12-2，书内 p.889—890 / PDF p.375—376 改绘。}}
\label{fig:stm-junction}
\end{figure}

\subsubsection{隧穿电流}

在解释图像衬度起源前，先要理解针尖与表面之间的电子流动过程\src{S2 §12.1.2，
书内 p.889—892 / PDF p.375—378}。两个金属电极靠近但未连通时，真空能级一致、各自费米
能级保持本征值，其差（接触电势）等于两金属功函数之差；用导线连通后真空能级发生偏移、
两者费米能级拉平，真空中出现电场，但费米能级齐平、没有可用能量窗口，故无净电流。只有
引入偏压才能打开能量窗口：被施加正偏压的一侧费米能级相对拉低 $eU$，其最高占据态与另一侧
的未占据态在能量上对齐，电子才发生隧穿；偏压反向则隧穿方向反转。两极相距较远时电流无法
穿越真空势垒，充分接近至 \qty{1}{nm} 左右时两侧电子波函数发生重叠，即出现隧穿电流。

\begin{principle}
取真空能级为能量参考点，金属表面功函数为 $\Phi$，则 $E_{\mathrm{F}}=-\Phi$。在针尖与
样品功函数相同的简化情形下，处于 $E_{\mathrm{F}}-eU$ 与 $E_{\mathrm{F}}$ 之间、能量为
$E_n$ 的样品态 $\varphi_n$ 有机会隧穿入针尖。当偏压远小于功函数值（$eU\ll\Phi$）时，
对电流有贡献的样品态主要位于费米能级附近，电子的隧穿可用恒定势垒下的薛定谔方程描述，
经典禁阻区内 $z$ 点附近的电子波函数为

\begin{equation}
\varphi(z)=\varphi(0)\,\mathrm{e}^{-\kappa z},
\qquad
\kappa=\frac{\sqrt{2m\Phi}}{\hbar}
\label{eq:kappa}
\end{equation}

式中 $\kappa$ 称衰减常数，$\varphi(0)$ 为样品表面处的电子态。于是第 $n$ 个样品态中的
电子出现在针尖表面处（$z=W$）的概率正比于
$\lvert\varphi_n(0)\rvert^{2}\mathrm{e}^{-2\kappa W}$。把能量窗口 $eU$ 内的所有样品态
都包括在内，隧穿电流为

\begin{equation}
I\;\propto\!\!\sum_{E_n=E_{\mathrm{F}}-eU}^{E_{\mathrm{F}}}
\lvert\varphi_n(0)\rvert^{2}\mathrm{e}^{-2\kappa W}.
\label{eq:sumI}
\end{equation}

按局域态密度（local density of states，LDOS）的定义，式~\eqref{eq:sumI} 可改写为

\begin{equation}
\begin{split}
I\;&\propto\;U\,\rho_{\mathrm{s}}(0,E_{\mathrm{F}})\,\mathrm{e}^{-2\kappa W}\\
   &\approx\;U\,\rho_{\mathrm{s}}(0,E_{\mathrm{F}})\,
   \mathrm{e}^{-1.025\sqrt{\Phi}\,W},
\end{split}
\label{eq:tunneling}
\end{equation}

其中 $\rho_{\mathrm{s}}(0,E_{\mathrm{F}})$ 为费米能级附近的表面局域态密度。注意
式~\eqref{eq:tunneling} 的指数上是 $2\kappa W$，故有实用单位定值的是 $2\kappa$ 而不是
$\kappa$：取 $m=\qty{9.109e-31}{kg}$、$\hbar=\qty{1.055e-34}{J.s}$，以 \unit{eV} 表示
$\Phi$、以 \unit{\angstrom} 表示 $W$，则
$2\kappa\approx\qty{1.025}{\per\angstrom}\times\sqrt{\Phi/\unit{eV}}$，即
$\kappa\approx\qty{0.512}{\per\angstrom}\times\sqrt{\Phi/\unit{eV}}$。系数归属已对照
原页面图像核对：原书该步印
$I\propto U\rho_{\mathrm{s}}(0,E_{\mathrm{F}})\mathrm{e}^{-2kW}
\approx U\rho_{\mathrm{s}}(0,E_{\mathrm{F}})\mathrm{e}^{-1.0252\sqrt{\varphi}W}$
（原书把这一衰减常数记作 $k$，本节记作 $\kappa$），$1.0252$ 确实位于整个指数上，即属于
$2\kappa$；本节按所列常数算得 $2\kappa\approx1.024$，与原书末位之差属常数取值与取位，
正文取 $1.025$ \src{S2 §12.1.2，书内 p.891 / PDF p.377}。
\end{principle}

式~\eqref{eq:tunneling} 给出两个必须分开记住的结论。

\textbf{（1）电流对间距指数敏感，带来原子级垂直分辨。} 取 $\Phi=\qty{4}{eV}$，由
式~\eqref{eq:kappa} 得 $\kappa\approx\qty{1.02}{\per\angstrom}$，即指数上的
$2\kappa\approx\qty{2.05}{\per\angstrom}$。这就是本书所说的\textbf{针尖与样品间距每增加
\qty{1}{\angstrom}，隧穿电流约衰减至 $\mathrm{e}^{-2}$（约 $1/7.4$）}；反过来，间距变化
\qty{0.1}{\angstrom} 对应约 23\% 的电流变化。若电流测量精度为 1\%，对应间距变化
$\Delta d=\ln(1.01)/(2\kappa)\approx0.0049$ \AA$\,=\qty{0.49}{pm}$，这正是 STM 纵向
灵敏度可达皮米量级的原因。
但\textbf{灵敏度不等于图像的实际纵向分辨率}：后者还受反馈环带宽、压电陶瓷迟滞与蠕变、
电子学噪声限制，恒流模式 $z$ 通道的实用分辨率通常在 \qtyrange{0.01}{0.1}{\angstrom} 量级。

\textbf{（2）测到的是电子态密度，不是纯几何高度。} 电流正比于
$\rho_{\mathrm{s}}(0,E_{\mathrm{F}})$，即能量窗口内样品表面电子态的数目，而非表面原子的
几何位置。被测的是哪批电子态由偏压极性决定：样品加正偏压时电子由针尖占据态流向样品空态，
探测 $E_{\mathrm{F}}$ 以上的未占据态；样品加负偏压时探测 $E_{\mathrm{F}}$ 以下的占据态。
因此调节偏压即可探测费米面附近不同的态密度，\textbf{同一表面在不同偏压下会给出不同的
像}。

这解释了一个看似矛盾的现象：Xe 吸附于 Ni(110) 时，吸附原子的电子态并不位于费米能级
附近，STM 却仍能很好地成像。原因是 Xe 显著改变了 Ni 的局域态密度——相比真空，Xe 降低了
局域的隧穿势垒，使表面波函数可以进一步延伸到真空中，从而被针尖在更远处探测到
\src{S2 §12.1.2，书内 p.891—892 / PDF p.377—378}。可见图像衬度取决于针尖与样品电子
结构的卷积以及两者间距：洁净金属表面上位置较高的原子通常对应凸起亮点，但对吸附质而言，
衬度还取决于其与基底的相互作用强度、偏压极性与大小、势垒高度的局域变化。相互作用较弱、
不显著改变表面电子结构的吸附物种仍可能被分辨出几何形态，但其表观高度不等于几何高度。

\begin{pitfall}
不要把 STM 图像当等高线图读。同一表面在不同偏压下会得到不同的像，一张像只对应一个能量
窗口内的电子态分布；把表观高度直接当作原子几何高度，是最常见也最难自查的误读。
\end{pitfall}

\subsubsection{恒高模式与恒流模式}

STM 成像有两种主要模式\src{S2 §12.1.1，书内 p.888 / PDF p.374}：\textbf{恒定电流模式}
对压电陶瓷扫描头施加反馈以保持隧穿电流恒定，同时记录为维持恒流所调整的 $z$ 位移随 $xy$
的变化；\textbf{恒定高度模式}保持针尖高度不变，直接测量隧穿电流随 $xy$ 的变化。两者的
差别不只是记录 $z$ 还是记录 $I$，而是\textbf{把指数敏感的那一端接在反馈环上，还是直接
送到信号通道}。

\begin{figure}[htbp]
\centering
\begin{tikzpicture}[
  x=1cm, y=1cm,
  base/.style={fill=black!12, draw=none},
  surf/.style={mainline},
  atom/.style={fill=black!30, draw=black, line width=0.4pt},
  ads/.style={fill=black!65, draw=black, line width=0.4pt},
  guide/.style={thinline, black, dashed},
  path/.style={mainline, black, dashed},
  ann/.style={font=\fontsize{6.6}{8}\selectfont, align=center},
  annl/.style={ann, align=left}
]
% ================= (a) 恒高模式 =================
\begin{scope}[shift={(0,0)}]
  \fill[base] (0.35,-0.05) rectangle (2.80,0.45);
  \fill[base] (2.80,-0.05) rectangle (4.40,0.62);
  \draw[surf] (0.35,0.45) -- (2.80,0.45) -- (2.80,0.62) -- (4.40,0.62);
  \foreach \xx in {0.62,0.95,1.28,1.61,1.94,2.27,2.60,3.05,3.45,3.85,4.20} {
    \fill[atom] (\xx,0.62) circle (1.1pt);
  }
  \fill[ads] (2.10,0.92) circle (1.5pt);
  \draw[thinline] (2.80,0.62) -- (2.80,0.45);
  \node[ann, align=left, anchor=west] at (2.62,0.02) {台阶};
  \fill[black!22] (1.35,2.30) -- (1.45,1.50) .. controls (1.70,1.36) and
                   (1.85,1.36) .. (2.10,1.50) -- (2.20,2.30) -- cycle;
  \draw[mainline] (1.35,2.30) -- (1.45,1.50) .. controls (1.70,1.36) and
                  (1.85,1.36) .. (2.10,1.50) -- (2.20,2.30);
  \draw[path] (0.40,1.55) -- (4.35,1.55);
  \draw[-{Stealth[length=2mm]}, thinline] (3.22,1.72) -- (4.16,1.72);
  \node[ann, anchor=south] at (3.88,1.76) {扫描方向};
  \draw[guide] (2.03,1.55) -- (2.03,0.99);
  \node[ann, align=left, anchor=west] at (1.05,1.90) {$z$ 预设恒定；\\ 读 $I(x,y)$};
  \node[ann, anchor=north] at (2.20,0.90) {$\Delta z$};
\end{scope}
% ================= (b) 恒流模式 =================
\begin{scope}[shift={(4.68,0)}]
  \fill[base] (0.35,-0.05) rectangle (2.80,0.45);
  \fill[base] (2.80,-0.05) rectangle (4.40,0.62);
  \draw[surf] (0.35,0.45) -- (2.80,0.45) -- (2.80,0.62) -- (4.40,0.62);
  \foreach \xx in {0.62,0.95,1.28,1.61,1.94,2.27,2.60,3.05,3.45,3.85,4.20} {
    \fill[atom] (\xx,0.62) circle (1.1pt);
  }
  \fill[ads] (2.10,0.92) circle (1.5pt);
  \draw[thinline] (2.80,0.62) -- (2.80,0.45);
  \fill[black!22] (1.35,2.30) -- (1.45,1.50) .. controls (1.70,1.36) and
                   (1.85,1.36) .. (2.10,1.50) -- (2.20,2.30) -- cycle;
  \draw[mainline] (1.35,2.30) -- (1.45,1.50) .. controls (1.70,1.36) and
                  (1.85,1.36) .. (2.10,1.50) -- (2.20,2.30);
  \draw[path]
    (0.40,1.05) -- (1.80,1.05) .. controls (1.95,1.05) and (1.98,1.35) ..
    (2.10,1.35) .. controls (2.24,1.35) and (2.28,1.05) .. (2.45,1.05) --
    (4.22,1.05);
  \draw[-{Stealth[length=2mm]}, thinline] (3.20,1.22) -- (4.16,1.22);
  % 【重叠修正】原为 anchor=south at (3.70,1.26)：文字向上展开到 y≈1.36，
  % 而下一行「Δz 被记录」在 (3.12,1.44) 向西对齐、x 覆盖 3.12—4.11，
  % 两者 x 区间几乎完全重合 —— 实测 p.114 出现「扫/Δ」「描/z」两对字符叠印。
  % 改为放到箭头下方（anchor=north），把上方让给 Δz 标注。
  \node[ann, anchor=north] at (3.70,1.16) {扫描方向};
  \node[ann, align=left, anchor=west] at (0.45,2.42) {$I$ 恒定；\\ 读 $z(x,y)$};
  \draw[guide] (2.72,1.30) -- (3.12,1.30);
  \node[ann, anchor=west] at (3.12,1.44) {$\Delta z$ 被记录};
\end{scope}
\end{tikzpicture}
\caption{恒高模式（a）与恒流模式（b）的原理对比。读者应看两点：针尖轨迹的高度是预设的
还是由反馈跟随表面决定的；信号通道记录的是电流还是 $z$ 位移。图中表面起伏与吸附质尺寸
为示意画法，未按真实比例。\schemamark\figsrc{据 S2（下册）§12.1.1 正文
（书内 p.888 / PDF p.374）重绘。}}
\label{fig:stm-modes}
\end{figure}

\begin{measure}
恒流模式的反馈环要求足够的响应速度：扫描越快、表面起伏越陡，环的滞后越明显，针尖轨迹会
削峰填谷，把形貌信息混入电子结构信息。恒高模式没有反馈环，扫描速率只受数据采集与前级
放大器带宽限制，但针尖与样品的最小间距在扫描区内无法自适应，因此要求该区粗糙度远小于
预设间距，偏压与电流设定值也须更保守以避免碰撞\src{S2 §12.1.1，书内 p.888 /
PDF p.374；设定值取舍为据原理的推断}。两者的适用场合不同：恒高模式不必等待反馈稳定，
常用于平整表面的原子分辨成像与快速扫描；恒流模式能安全跨越台阶与纳米颗粒，是粗糙表面、
原位与变温实验的默认选择。正因反馈电路与针尖稳定性的限制，目前的 STM 研究较多集中于
平整表面，以获得原子尺度的空间分辨成像\src{S2 引言，书内 p.887 / PDF p.373}。
\end{measure}

\subsubsection{扫描隧道谱（STS）}

STM 的能力不止于形貌。式~\eqref{eq:tunneling} 已表明图像不仅取决于表面形貌，也取决于加在
针尖上的电压——对电压的控制决定了哪个电子态发生隧穿。因此 STM 可用于原子尺度的局域电子
结构测量，这种数据采集模式称为扫描隧道谱（scanning tunneling spectroscopy，STS）
\src{S2 §12.1.3，书内 p.892 / PDF p.378}。

\begin{principle}
按 WKB 近似，隧穿电流可写成显式积分形式：
\begin{equation}
I\;\propto\!\int_{0}^{eU}\!
\rho_{\mathrm{s}}(\mathbf{r},E)\,
\rho_{\mathrm{t}}(\mathbf{r},\pm eU\mp E)\,
T(E,eU,\mathbf{r})\,\mathrm{d}E ,
\label{eq:wkb}
\end{equation}
式中 $\rho_{\mathrm{s}}(\mathbf{r},E)$ 与 $\rho_{\mathrm{t}}(\mathbf{r},E)$ 分别为样品与
针尖在位置 $\mathbf{r}$、能级 $E$ 处的电子态密度；$+eU-E$ 对应样品加正偏压，
$-eU+E$ 对应样品加负偏压；$T$ 为隧穿概率函数。取针尖态密度在讨论的能量范围内近似平坦，
并对偏压求导——低温极限下费米分布的导数趋于 $\delta$ 函数——于是
\begin{equation}
\left.\frac{\mathrm{d}I}{\mathrm{d}V}\right|_{\mathbf{r}}
\;\propto\;\rho_{\mathrm{s}}(\mathbf{r},E_{\mathrm{F}}+eV)\;
T(E_{\mathrm{F}},eV,\mathbf{r}).
\label{eq:dIdV}
\end{equation}
即\textbf{微分电导正比于针尖所在位置的样品局域态密度}。这就是 STS 测的是电子结构而非
形貌的原因：式~\eqref{eq:dIdV} 的右端只含态密度与隧穿概率，几何高度没有独立地位。
式~\eqref{eq:wkb} 的积分限与能量宗量已逐字核对：原书印
$I=\int_{0}^{eU}\rho_{\mathrm{s}}(r,E)\rho_{\mathrm{t}}(r,\pm eU\mp E)T(E,eU,r)\mathrm{d}E$，
并写明「$+eU-E$ 对应对样品施加正偏压 $eU$ 的情况；$-eU+E$ 对应于对样品施加负偏压 $eU$
的情况」（原书用等号，本节省去前置常数故写正比号）。原书紧接的一句是「归一化后的微分电导
将只取决于局域态密度和隧穿电压」，确以 $I/V$ 归一化为前提；式~\eqref{eq:dIdV} 是本节据该
表述补出的求导形式，原书未列出式~\eqref{eq:dIdV}。
\src{S2 §12.1.3，书内 p.892 / PDF p.378，STS 一节的隧穿电流积分式}
\end{principle}

\begin{dataflow}
STS 的数据链条是：针尖固定在 $\mathbf{r}$ 处扫偏压得到 $I(V)$ 曲线（非线性）$\rightarrow$
取微分电导 $\mathrm{d}I/\mathrm{d}V$ $\rightarrow$ 按式~\eqref{eq:dIdV} 换算为费米能级附近
的局域态密度 $\rho_{\mathrm{s}}(\mathbf{r},E_{\mathrm{F}}+eV)$ $\rightarrow$ 在 $xy$ 上
逐点重复得到 LDOS 空间成像。实际操作中常在直流偏压上叠加小幅高频调制，用锁相放大器取
电流的一阶谐波直接得到 $\mathrm{d}I/\mathrm{d}V$（这是 STS 的通行做法，非本书记载）。
STS 之所以必要，是因为它能把几何上相似的对象分开：氧化物表面的桥氧空位与轻原子吸附物种
在形貌图上可能都表现为亮点，而同一次实验采集的 $\mathrm{d}I/\mathrm{d}V$ 曲线或映射图可按
电子结构把它们分辨开；反过来，表观高度不同的两个物种在 STS 上也可能给出相同的谱，此时
不能据形貌图推断活性位不同。
\end{dataflow}

STS 在纳米催化中的一个典型用法是测量局域功函数与局域态密度随结构参数的变化，见
\ref{sec:model-nano} 节。六环节在本节的落点是：\textbf{问题与输出}确定 STM 给出 $I(x,y)$
或 $z(x,y)$；\textbf{原理与公式}（式~\eqref{eq:tunneling}）说明电流为何对间距指数敏感、
又为何编码态密度；\textbf{测量与条件}给出恒高与恒流模式的取舍；\textbf{数据到结论}把
$z$ 通道换算为形貌、把 $\mathrm{d}I/\mathrm{d}V$ 换算为局域态密度；\textbf{误判与边界}
与\textbf{互补与回查}分别见后两节。

\subsection{STM 在催化研究中的应用}

\subsubsection{催化反应机理和动力学的可视化研究}

随着低温和变温 STM 技术的发展，STM 可以实时监测化学反应，即反应物的化学吸附、扩散和催化
转化的基元步骤\src{S2 §12.2.1，书内 p.892—903 / PDF p.378—389}。低温下表面扩散被抑制、
吸附物种的运动完全可控，有利于中间体鉴定；STM 还可用针尖脉冲诱发吸附、脱附与解离。在
Ag(110) 上，针尖在分子氧顶端施加 \qty{0.47}{V} 脉冲使其解离为一对氧原子，\qty{0.24}{V}
脉冲可驱动 CO 迁移；CO 靠近氧原子对时自动形成 O—CO—O 复合体，复合体在脉冲下解离，
生成的 \ce{CO2} 脱附而表面留下一个氧原子，全程遵循 Langmuir--Hinshelwood 机制。针尖操纵
与成像结合，说明 STM 能在单分子水平做反应路径成像。

\textbf{从扩散看活性位。} 覆盖度增加时吸附分子间的相互吸引导致吸附质岛形成，并与自由移动
的分子保持动态平衡。跟踪表面空位扩散可得到活性位信息：在约 \qty{65}{K} 与
\qty{2e-7}{torr}（\qty{1}{torr}=\qty{133.32}{Pa}）的 \ce{H2} 中 Pd(111) 表面可被 H 饱和，
由于针尖吸附了 H 原子，成像机制反转，\textbf{吸附空位表现为凸起亮点}。连续图像显示空位
二聚体会重新解离为孤立空位，而\textbf{三聚体消失只留下一个空位}，说明其中两个空位被
\ce{H2} 解离产生的两个 H 占据；大量统计表明只有三个或更多空位的聚集体才会被 H 占据，即
\ce{H2} 在 Pd 表面的解离需要三个或更多紧邻空位\src{S2 §12.2.1，书内 p.894—895 /
PDF p.380—381}——这与「两个相邻空位已足够」的传统默认不同。

\textbf{氧化物表面。} 金红石 \ce{TiO2}(110) 表面桥氧加热时部分脱附形成桥氧空位，这是该表面
最常见的缺陷位。结合原位 STM 与 DFT 可把活性位、吸附物种与中间体区分开：常用正偏压扫描下
5 配位 Ti 原子列成像为亮行、桥氧列成像为暗行，桥氧空位、表面 OH 与吸附 \ce{H2O} 有明显
的表观高度差异。由此得到：\ce{H2O} 优先在桥氧空位吸附解离，解离的 H 落到近邻桥氧上形成
配对 OH；测量 300—\qty{410}{K} 的迁移速率得，相邻桥氧位上的 H（$H_{\mathrm{B}}$）的迁移
能垒比空位 OH 上的 H（$H_{\mathrm{V}}$）低约 \qty{0.22}{eV}，且能垒随 OH 间距增大而增加，
表明 OH—OH 排斥帮助了 H 的迁移；甲醇在桥氧空位解离吸附形成的甲氧基比桥氧空位高
\qty{0.8}{\angstrom}，随后 OH 中 H 的迁移可被逐帧跟踪\src{S2 §12.2.1，书内 p.895—899 /
PDF p.381—385}。

\textbf{反应机理。} 第 \textsc{viii} 族贵金属催化的氢氧化反应是 1823 年发现的第一个催化
反应，其机理长期不清楚：通常认为经 \ce{O_{ad}} 与 \ce{H_{ad}} 结合成羟基、再与另一个
\ce{H_{ad}} 成水，羟基形成常被假定为决速步，但羟基难以在早期光谱研究中证实。在
\ce{Pt}(111) 上预吸附 \ce{O2} 形成 $(2\times2)$--\ce{O_{ad}}（\ce{O_{ad}} 为暗点）后通入
\ce{H2}：上方出现亮岛并扩大，最终形成氢键连接的 \ce{OH_{ad}} 有序覆盖层，图像中还可见
迁移的小团簇（\ce{H2O_{ad}}）。大尺度上可见随时间向外扩展的亮环，即\textbf{反应前沿}：
边界处 \ce{O_{ad}} 与扩散 \ce{H2O_{ad}} 快速反应生成两个 \ce{OH_{ad}}，后者再与
\ce{H_{ad}} 快速成水。由于 \ce{O_{ad}}+\ce{H_{ad}} 成羟基是决速步，\ce{H2O_{ad}} 的参与
消除了这一限制，形成\textbf{自催化}过程，这解释了氢氧化在水的脱附温度以下仍有高活性的
来源；当 \ce{H2O} 开始脱附、\ce{H2O_{ad}} 寿命缩短，自催化循环即被打破
\src{S2 §12.2.1，书内 p.900—901 / PDF p.386—387}。

\begin{dataflow}
\textbf{反应动力学。} 单晶表面 CO 氧化的原位 STM 研究通常用\textbf{滴定实验}：预吸附
\ce{O_{ad}} 后暴露于恒定压力的 CO，速率直接由覆盖度随时间的变化得到。在 \ce{Pt}(111) 上，
\qty{247}{K} 下 \ce{O_{ad}} 形成有序 $(2\times2)$--O（暗点）；通入 \qty{5e-8}{torr} CO 后
CO 形成 $c(4\times2)$ 并扩大，$(2\times2)$--O 岛面积减小。把面积减小速率对周长（两类区域
的畴界长度）与对 $\theta_{\mathrm{O}}(1-\theta_{\mathrm{O}})$ 分别作图，可见\textbf{速率
与周长线性相关}，说明 CO 氧化主要沿边界发生；在 237—\qty{274}{K} 重复滴定，求得活化能约
\qty{0.49}{eV}、指前因子约 \qty{3e21}{cm^{-2}.s^{-1}}，与宏观测量相符
\src{S2 §12.2.1，书内 p.901—903 / PDF p.387—389}。并非所有体系都能用简单级数描述：
\ce{Ag}(110) 上 \ce{O2} 解离引起的 $(n\times1)$--O 重构由一维 AgO 链构成，室温暴露于
\qty{1e-8}{torr} CO 时链分段消失、速率显著加快；低于 \qty{230}{K} 时反应只发生在链末端，
为零级动力学，能垒约 \qty{41}{kJ/mol}。
\end{dataflow}
\subsubsection{近常压 STM 研究}

上节的研究暴露出一个根本疑虑：化学吸附可改变催化剂表面结构或组分，超高真空下的结果能否
外推？此前表征常在 \qtyrange{1e-14}{1e-10}{bar} 的 UHV 下进行，而工业催化在
\qtyrange{1}{1000}{bar} 下进行，超过十个数量级的压差可极大改变反应物在表面的相互作用。
要跨越这一\textbf{压力鸿沟}，原位 STM 的观测必须拓展到近常压甚至高压
\src{S2 §12.2.2，书内 p.903—904 / PDF p.389—390}。Salmeron 等较早把常压 STM 用于模型
催化剂，展示了 \ce{Pt}(110) 在接近一个大气压的 \ce{H2}、\ce{O2} 或 CO 中的重构。Frenken
等进一步把扫描头置于流动反应器内：UHV 中 \ce{Pt}(110) 重构为 $(1\times2)$，CO 可消除该
重构；在 \qty{1}{bar} CO 下高压 CO 不仅消除重构还使表面粗糙化——\qty{425}{K}、
\qty{1.25}{bar} CO 下由 $(1\times2)$ 向 $(1\times1)$ 的重排导致表面碎片化并形成高密度
台阶（$t=0$ min 呈虎皮状），随后台阶熟化、表面平滑并形成 CO 覆盖的大岛。

把 CO/\ce{O2} 混合气通入反应器，同时用四极质谱仪实时监测 CO、\ce{O2} 与 \ce{CO2}，即可
实现表面结构与反应性能的同时测量，即 \textbf{operando} STM。按速率可分为两阶段：低活性
阶段图像无明显变化，Pt 保持金属态；高活性阶段速率约为前者的 3 倍，表面粗糙、可能有铂
氧化物形成，且粗糙度随速率均匀增加。结合高压 SXRD（该相被指认为铂氧化物单层）与动力学
测量，Frenken 等提出高活性阶段 CO 可直接与表面氧化物的氧反应，即 Mars--van Krevelen
机制；但书中明确指出，由于高压环境的复杂性（如气体输运受限时表面附近气相可与远处气相
大不相同），该构效关系仍存争议，还需结合谱学做更多 operando 研究
\src{S2 §12.2.2，书内 p.905—906 / PDF p.391—392}。高压 STM 的必要性正在于此：高压气体
造成的吸附物种结构常可在高真空低温实验中重现，但金属原子在低温下运动受限，表面结构在
实际催化条件下的\textbf{动态}变化难以在 UHV 中直接观察。
\subsubsection{负载型纳米结构催化剂的模型研究}
\label{sec:model-nano}

实际催化剂常由 \qtyrange{1}{10}{nm} 的活性纳米颗粒分散在高比表面载体上构成，颗粒可具有与
体相完全不同的结构和性质，由此产生与单晶平整表面之间的\textbf{材料鸿沟}。反应性和选择性
常取决于颗粒的大小、形貌及与载体的相互作用，即\textbf{尺寸效应}与\textbf{界面效应}，而
这两个效应无法从平整单晶表面的研究得到\src{S2 §12.2.3，书内 p.907 / PDF p.393}。

\textbf{金属-氧化物界面位点。} 包信和团队在 \ce{Pt}(111) 表面构筑 FeO--Pt 界面位点，发现
Pt 表面可稳定配位不饱和的 \ce{Fe^{2+}}（coordinatively unsaturated ferrous，CUF）中心；
该位点在室温 Pt 表面被 CO 毒化时仍保持很高的 \ce{O2} 活化能力，可实现富氢条件下 CO 的
选择性氧化（PROX）。进一步构筑不同尺寸的 FeO 纳米三角形岛后，由 STM 图定出各岛的化学
计量比（如 \ce{Fe66O55}、\ce{Fe190O171}、\ce{Fe378O351}），边界均呈现相同的 CUF 结构。
STS 结果显示：局域功函数（local work function，LWF）与局域态密度（LDOS）随 FeO 在
\ce{Pt}(111) 上的\textbf{落位}而变——Fe 落在 fcc 位点时 LWF 最大（\qty{5.1}{eV}），落在
top 位点时最低（\qty{4.2}{eV}）；但边长 $>\qty{3}{nm}$ 的不同尺寸纳米岛在相同落位（即
相同 Pt--Fe--O 界面结构）表现出相同的 LWF 与 LDOS，不受尺寸影响。LDOS 空间成像还显示，
边界靠近 top 位点及 top-FeO 区域在费米能级以上约 \qty{0.5}{eV} 处呈共振峰
\src{S2 §12.2.3，书内 p.907—908 / PDF p.393—394}。

这里出现一个值得记录的反差。传统认识上相同的活性结构与电子结构必然带来相同的活性，但用
原位 STM 观察不同尺寸 FeO 纳米岛与 \ce{O2} 的相互作用时，\ce{O2} 在 CUF 边界的吸附解离
同时带来岛的\textbf{动态重构}，且明显依赖尺寸：边长约 \qty{3}{nm} 的岛在氧解离后整个岛面
Fe--O 位置相对偏移，最终整个岛的氧晶格重构为稳定的二配位氧终止边；而边长大于
\qty{4}{nm} 的岛只能部分重构，未重构的 CUF 边界吸附了相对不稳定的氧原子，岛面出现重构区
与未重构区之间的畴界。这种与尺寸相关的氧气氛动态重构被称为\textbf{动态尺寸效应}，显著
影响 FeO 纳米粒子在 \ce{O2} 中的稳定性：在 \qty{500}{K}、\qty{1e-5}{mbar} \ce{O2} 中处
理时，尺寸小于 \qty{3.2}{nm} 的岛因更易完整重构、边界形成惰性的二配位氧终止边而表现出
反常的抗氧化行为，大岛则因重构不完全、边界存留的不稳定氧较易从 CUF 边或畴界扩散进入而
深度氧化\src{S2 §12.2.3，书内 p.908—910 / PDF p.394—396}。

\begin{infer}
上例说明：STS 显示的「电子结构不随尺寸变化」与原位 STM 显示的「反应行为随尺寸变化」并不
矛盾——前者描述反应前特定落位上的电子结构，后者描述\textbf{吸附诱导重构后}的终态。若只用
STS 或形貌图之一推断活性，会把反应过程中的结构演化漏掉。
\end{infer}

STM 也能解析实际纳米催化剂结构。Goodman 等在 Mo(112) 上预制备 \ce{SiO2} 单层薄膜，把分子
团簇催化剂 \ce{Ru3(CO)9(SnPh2)3}（Ph 为苯基）经溶液分散沉积到膜上并在真空中加热去除 CO
与苯基配体；高分辨 STM 图显示其原子分辨图包含\textbf{六个金属原子}，排列与前体团簇中的
\ce{Ru3Sn3} 相同，表明配体低温去除时团簇骨架在氧化硅表面保持完整
\src{S2 §12.2.3，书内 p.910 / PDF p.396}。

\textbf{金属载体强相互作用（SMSI）。} 20 世纪 70 年代 Tauster 等发现第 \textsc{viii} 族
过渡金属（如 Pt、Ir）负载在 \ce{TiO2} 上并在较高温度还原时表现出不寻常特性（对 CO 和
\ce{H2} 的化学吸附受抑制而甲烷化活性增加），据此提出 SMSI 概念，其起源至今仍是热点。
Bowker 等把平均约 \qty{4}{nm} 的 Pd 颗粒负载于 \ce{TiO2}(110)，在 \qty{673}{K}、
\qty{5e-8}{mbar} \ce{O2} 下做原位 STM：\ce{O2} 在 Pd 上的解离效率比在 \ce{TiO2}(110) 上
高 3 个数量级，故解离主要发生在 Pd 上，解离的 O 随即溢流到相邻表面并与扩散到表面的间隙
态 \ce{Ti^{3+}} 结合，使 Pd 周围 \ce{TiOx} 富集；\ce{TiOx} 层再生长加速并最终超过 Pd
颗粒高度将其包裹。原位 STM 因此证实\textbf{氧化物包裹}是 SMSI 的重要原因
\src{S2 §12.2.3，书内 p.910—911 / PDF p.396—397}。

\textbf{Au 纳米颗粒的烧结动力学。} 纳米粒子的烧结有两种传质模式：整个粒子的迁移与团聚，
以及单体（单个金属原子或配合物）的迁移与颗粒熟化（Ostwald 熟化）。Goodman 等在 UHV 与近
常压条件下研究 Au/\ce{TiO2}(110)：真空下 Au 颗粒在室温保持稳定，而引入 \qty{0.1}{torr}
的 CO/\ce{O2} 混合气（$1\!:\!1$）后立刻开始烧结。连续原位 STM 显示粒子\textbf{位置没有
变化}，而小粒子逐渐变小消失、大粒子长大，直接表明发生了 Ostwald 熟化而非粒子迁移；单独
暴露于 \ce{O2} 或 CO 时几乎无变化，故推断是混合气的协同作用诱导并加速了烧结。基于
Ostwald 熟化模型，对直径 \qty{3}{nm} 的 Au 粒子求得表观熟化活化能约 \qty{10}{kJ/mol}，
与该催化剂上 CO 氧化的活化能相近\src{S2 §12.2.3，书内 p.911—912 / PDF p.397—398}。
\subsection{误判与边界}
\label{sec:stm-pitfall}

\subsubsection{针尖卷积与「双针尖」假象}

STM 图像是针尖形状与样品表面电子结构的\textbf{卷积}，不是表面的直接投影。理想针尖末端只有
一个原子突起，实际针尖常有多个突起：每个突起各形成一条隧穿通道，图像便是这几条通道的
叠加，表现为\textbf{重影、孪生特征或特征被复制到邻近位置}，即「双针尖」（double tip）
假象。判据是：重影的出现与消失依赖于针尖处理，改变扫描方向时重影的相对位置发生变化，
同一样品更换针尖后重影消失。针尖卷积还会放大小尺寸特征的表观尺寸、抹圆单原子台阶、使
孤立吸附质的表观高度偏离几何高度，极端情形下把同一吸附质成像两次而使覆盖度被高估。另一
类反转与针尖状态有关：Pd(111) 上吸附了 H 原子的针尖会把\textbf{吸附空位}成像为凸起亮点
\src{S2 §12.2.1，书内 p.894 / PDF p.380}，此时按几何直觉解读会把空位误认为吸附物种、把
空位团簇的消失误认为吸附物种脱附；\ce{TiO2}(110) 上的氧空位同样成亮点是反例。

\begin{pitfall}
凡要从 STM 图读出「有几个原子」「覆盖度是多少」「哪个是吸附质」的结论，都必须先排除针尖
卷积：用不同针尖重复、改变扫描方向、用针尖脉冲处理后再成像，三种办法至少用两种。只凭
一张像下的结论不可靠。
\end{pitfall}

\subsubsection{STM 图像是电子态密度图}

同一表面在改变偏压时会给出不同的像，这是原理决定的，不是仪器故障——原因见
式~\eqref{eq:tunneling} 与式~\eqref{eq:dIdV}：不同偏压对应不同能量窗口。因此「STM 图上的
亮点」必须与偏压一起报告；不同文献中同一体系的表观高度不一致，往往源于偏压符号与数值不同
或针尖末端态密度差异，而不一定是样品不同。由此派生三条常见误读：把表观高度差直接当作
几何高度差；把亮点一律当作位置更高的原子（Pd(111) 的 H 空位、\ce{TiO2}(110) 的空位都是
反例）；把不同偏压下的两张像逐点相减当作同一批原子的位移。

\subsubsection{近常压 STM 中针尖与气体的相互作用}

近常压环境本身会改变测量。针尖在气体中可能被 CO 或其他物种覆盖，改变其电子结构与有效
功函数，甚至如 Pd/H 的例子那样反转成像机制；针尖-样品间隙中的气体分子也可能参与隧穿或
改变势垒高度，使电流-电压关系偏离式~\eqref{eq:tunneling} 的简化形式；高气压下气体输运
受限，表面附近气相成分可与远处气相不同，因此把质谱测得的体相压力直接当作表面反应条件是
不安全的\src{S2 §12.2.2，书内 p.905—906 / PDF p.391—392}。此外，operando 条件下扫描速率
与反应速率可能处于同一时间尺度，图像中看到的「结构」有可能是扫描期间的时平均结果。

\subsubsection{模型体系结论外推的边界}

这一条决定 STM 结果在催化中的适用范围。前述 12.2 的绝大多数结论来自\textbf{超高真空中、
单晶（或单晶上的有序氧化物纳米岛）表面、低温或中等温度}下的模型体系。外推到真实负载型
催化剂时至少有三重边界\src{S2 §12.2.2—12.2.3，书内 p.903—912 / PDF p.389—398}。

\textbf{（1）压力边界。} UHV 表征常在 \qtyrange{1e-14}{1e-10}{bar}，工业催化在
\qtyrange{1}{1000}{bar}，相差十个数量级以上。高压下表面可发生 UHV 中不出现的重构（如
\ce{Pt}(110) 在 \qty{1}{bar} CO 下的粗糙化与台阶熟化）甚至形成新的表面相（高活性阶段的
铂氧化物）。低温高覆盖度实验能复现高压下的吸附结构，却不能复现高压下的表面动力学，因为
金属原子在低温下运动受限。

\textbf{（2）材料边界（材料鸿沟）。} 真实催化剂是 \qtyrange{1}{10}{nm} 的颗粒分散在高比
表面载体上，颗粒有尺寸分布、多种暴露晶面与载体接触界面；模型体系则是大面积单晶平台或
尺寸单一、边界结构统一的有序纳米岛，尺寸效应与界面效应恰好在其中被标准化掉了。FeO 纳米
岛的例子显示了这种标准化的代价：边长大于 \qty{3}{nm} 的岛在相同落位上电子结构相同，但
一旦与 \ce{O2} 作用，重构行为就重新表现出尺寸依赖性，而这个依赖性在静息的电子结构上看
不出来。模型体系的有序度本身也是精心制备的产物，真实载体上的分散状态与之差距很大。

\textbf{（3）方法边界。} STM 只能测导电或半导体样品，只能看到表面最外层的电子态分布，
本身\textbf{不提供化学识别}——区分物种要靠 STS、针尖操纵脱附实验或与谱学手段回查，因此
从 STM 图出发的活性位归属在原理上是间接结论。

综合起来，STM 结论可靠的外推范围是「与模型体系在同量级压力、同量级尺寸、同一表面化学
环境」。跨越压力鸿沟或材料鸿沟的推论必须重新取证：要么在近常压/高压下做 operando 测量，
要么在真实催化剂（或其高分散模型）上重复，要么由谱学与理论计算独立支持。

\subsection{互补与回查}
\label{sec:stm-cross}

形貌图与同一区域的 STS 映射应一起解读：前者给出几何形貌，后者给出费米面附近的局域态密度；
若两者对活性位的归属矛盾，应优先怀疑针尖状态或偏压选择，而不是样品本身。STS 测得的局域
功函数须与独立手段对照——宏观功函数测量、Kelvin 探针、紫外光电子能谱（UPS）给出面平均或
较大面积的平均值，STS 给出单个落位上的值，二者不一致是正常的（前者是平均、后者是分布），
但不能矛盾到量级不符；FeO/\ce{Pt}(111) 上 fcc 落位 \qty{5.1}{eV} 与 top 落位
\qty{4.2}{eV} 的差异需要 DFT 独立支持才能作为「落位决定电子结构」的结论。

动力学参数必须与宏观动力学对照。\ce{Pt}(111) 上 CO 氧化正是一个成功回查：由 STM 滴定求得
的活化能约 \qty{0.49}{eV} 与指前因子约 \qty{3e21}{cm^{-2}.s^{-1}} 与宏观测量相符，这既
验证了 STM 测速的可靠性，也说明由「速率正比于畴界长度」推出的「反应沿边界发生」这一微观
图像是自洽的。高压/operando 结构必须与高压谱学对照：Frenken 等对高活性阶段表面相的归属
依赖高压 SXRD 的独立结果，即便如此书中仍强调该构效关系存在争议，说明\textbf{单一
operando STM 不足以定论}。最后是模型体系与真实催化剂的回查：从单晶或有序纳米岛得到的
尺寸效应、界面效应结论，都应在真实负载催化剂或高分散模型催化剂上检验其边界。

\begin{crosscheck}
STM 提供空间分辨，STS 提供局域电子结构，二者仍不提供化学识别。判断物种归属时至少需要两条
独立证据链：一条来自图像（表观高度、针尖操纵脱附、偏压依赖性），另一条来自谱学或理论
计算。缺少第二条链时，把「亮点」直接命名为某个吸附物种是超出数据支持的推断。
\end{crosscheck}
